\documentclass[11pt,a4paper]{article}

\usepackage{fontspec}
\IfFontExistsTF{Times New Roman}{\setmainfont{Times New Roman}}{\setmainfont{Tinos}}
\IfFontExistsTF{Consolas}{\setmonofont{Consolas}[Scale=0.88]}{\setmonofont{Liberation Mono}[Scale=0.88]}
\usepackage[english]{babel}
\usepackage{geometry}
\usepackage{xcolor}
\usepackage{enumitem}
\usepackage{listings}
\usepackage{booktabs}
\usepackage{array}
\usepackage{amsmath,amssymb}
\usepackage{titlesec}
\usepackage{fancyhdr}
\usepackage{parskip}
\usepackage{tcolorbox}
\tcbuselibrary{skins}
\usepackage{tabularx}
\usepackage{lastpage}
\usepackage{hyperref}
\usepackage{multicol}

\geometry{top=1.3cm,bottom=1.3cm,left=1.65cm,right=1.65cm}

\definecolor{primaryblue}{RGB}{18,65,115}
\definecolor{secondaryteal}{RGB}{14,102,110}
\definecolor{SoftBlue}{RGB}{232,242,250}
\definecolor{CodeBg}{RGB}{248,250,252}
\definecolor{DarkGray}{RGB}{55,55,55}
\definecolor{boxbg}{RGB}{248,250,252}
\definecolor{borderblue}{RGB}{150,185,215}

\setlength{\parindent}{0pt}
\setlength{\parskip}{0pt}
\setlist[itemize]{leftmargin=1.2em,itemsep=0.3pt,topsep=0.3pt,parsep=0pt,partopsep=0pt}
\setlist[enumerate]{leftmargin=1.3em,itemsep=0.5pt,topsep=0.3pt,parsep=0pt,partopsep=0pt}
\setlist[enumerate,2]{topsep=0.1pt,itemsep=0pt,parsep=0pt,partopsep=0pt}

\lstset{
  language=Python,
  basicstyle=\ttfamily\scriptsize,
  keywordstyle=\color{blue!75!black}\bfseries,
  stringstyle=\color{red!65!black},
  commentstyle=\color{primaryblue!80!black}\itshape,
  showstringspaces=false,
  columns=fullflexible,
  keepspaces=true,
  breaklines=true,
  frame=single,
  rulecolor=\color{primaryblue!45},
  backgroundcolor=\color{CodeBg},
  xleftmargin=0.15em,
  xrightmargin=0.15em,
  aboveskip=0.10em,
  belowskip=0.10em,
  tabsize=4
}

\pagestyle{fancy}
\fancyhf{}
\setlength{\headheight}{14pt}
\fancyhead[L]{\footnotesize\textbf{DSAI1003 -- Fundamental Programming Concepts in Python}}
\fancyhead[R]{\footnotesize\textbf{Midterm Exam (Loops \& Functions, 60 Min) --- Code: 202}}
\fancyfoot[L]{\scriptsize Faculty of Data Science and Artificial Intelligence -- National Economics University}
\fancyfoot[R]{\footnotesize Page \thepage\ of 6}
\renewcommand{\headrulewidth}{0.4pt}
\renewcommand{\footrulewidth}{0.3pt}

\titleformat{\section}{\normalsize\bfseries\color{primaryblue}}{}{0pt}{}
\titleformat{\subsection}{\small\bfseries\color{primaryblue}}{}{0pt}{}
\titlespacing*{\section}{0pt}{2.0pt}{0.5pt}
\titlespacing*{\subsection}{0pt}{1.5pt}{0.4pt}

\newtcolorbox{headerbox}{
  colback=white,
  colframe=primaryblue,
  boxrule=0.7pt,
  arc=1.2mm,
  left=2.0mm,
  right=2.0mm,
  top=0.8mm,
  bottom=0.8mm
}

\newtcolorbox{answerbox}[2][]{
  enhanced,
  colback=white,
  colframe=DarkGray!45,
  boxrule=0.5pt,
  arc=1mm,
  left=2mm,
  right=2mm,
  top=0.8mm,
  bottom=0.8mm,
  title=\footnotesize\color{DarkGray}\textbf{#2},
  #1
}

\hypersetup{
  colorlinks=true,
  linkcolor=primaryblue,
  urlcolor=primaryblue,
  pdftitle={DSAI1003 - Progress Midterm Examination Code 202 (Loops and Functions, 60 Minutes)},
  pdfauthor={Faculty of Data Science and Artificial Intelligence - National Economics University (NEU)}
}

\begin{document}

% =========================================================================
% PAGE 1: OFFICIAL HEADER + SCORE TABLE + ANSWER SHEET + Q1 - Q5
% =========================================================================

\noindent
\begin{minipage}[t]{0.64\textwidth}
  \raggedright\small
  \textbf{NATIONAL ECONOMICS UNIVERSITY}\\
  \textbf{COLLEGE OF TECHNOLOGY}\\
  \textbf{FACULTY OF DATA SCIENCE \& ARTIFICIAL INTELLIGENCE (FDA)}
\end{minipage}%
\hfill
\begin{minipage}[t]{0.34\textwidth}
  \raggedleft\footnotesize
  \textbf{OFFICIAL PROGRESS EXAMINATION}\\
  Academic Year: 2026--2027\\
  Semester: Fall 2026 \quad Course: \textbf{DSAI1003}
\end{minipage}

\vspace{0.08em}

\begin{center}
  {\large\bfseries\color{primaryblue} FUNDAMENTAL PROGRAMMING CONCEPTS IN PYTHON}\\[0.04em]
  {\normalsize\bfseries PROGRESS MIDTERM EXAM: LOOPS, ITERATION \& MODULAR FUNCTIONS}\\[0.04em]
  {\footnotesize\textbf{Paper Code:} \textbf{202} \quad | \quad \textbf{Duration: 60 Minutes} \quad | \quad \textbf{Difficulty:} 80\% Basic, 20\% Medium}
\end{center}

\vspace{0.04em}

\begin{headerbox}
\begin{tabular*}{\linewidth}{@{\extracolsep{\fill}}ll@{}}
  \textbf{Candidate Full Name:} \dotfill & \textbf{Student ID (MSSV):} \dotfill \\[0.10em]
  \textbf{Class / Cohort:} \dotfill & \textbf{Exam Room / Seat No.:} \dotfill
\end{tabular*}

\vspace{0.12em}
\renewcommand{\arraystretch}{1.0}
\setlength{\tabcolsep}{3.0pt}
\begin{tabularx}{\linewidth}{|l|c|c|c||c|X|}
\hline
\textbf{Section} & \textbf{\shortstack{Sec A: Quiz\\(40\%)}} & \textbf{\shortstack{Sec B: Fill-in\\(30\%)}} & \textbf{\shortstack{Sec C: Code\\(30\%)}} & \textbf{\shortstack{Total\\(100\%)}} & \centering\arraybackslash \textbf{\shortstack{Examiner\\Signature}} \\
\hline
\textbf{Max Score} & 40 pts & 30 pts & 30 pts & \textbf{100 pts} & \\[0.22em]
\cline{1-5}
\textbf{Score} & & & & & \\[0.32em]
\hline
\end{tabularx}

\vspace{0.10em}
\raggedright
\scriptsize
\textbf{Instructions:} Closed-book examination. All electronic devices and calculators are prohibited. Total score: 100 points (40\% Quiz, 30\% Fill-in, 30\% Full Code; 80\% Basic, 20\% Medium). Mark MCQs directly in the answer sheet grid. Write code solutions in designated boxes. Do not detach pages.
\end{headerbox}

\vspace{1.0mm}
\noindent
\textbf{MULTIPLE-CHOICE ANSWER SHEET} \textit{(Write letters A, B, C, or D into corresponding boxes)}:
\begin{center}
\setlength{\tabcolsep}{4.5pt}
\renewcommand{\arraystretch}{0.92}
\begin{tabular}{|c|c|c|c|c|c|c|c|c|c|c|}
\hline
\textbf{Q} & \textbf{1} & \textbf{2} & \textbf{3} & \textbf{4} & \textbf{5} & \textbf{6} & \textbf{7} & \textbf{8} & \textbf{9} & \textbf{10} \\
\hline
\textbf{Ans} & & & & & & & & & & \\[1.4mm]
\hline
\hline
\textbf{Q} & \textbf{11} & \textbf{12} & \textbf{13} & \textbf{14} & \textbf{15} & \textbf{16} & \textbf{17} & \textbf{18} & \textbf{19} & \textbf{20} \\
\hline
\textbf{Ans} & & & & & & & & & & \\[1.4mm]
\hline
\end{tabular}
\end{center}

\vspace{2.5mm}
\section*{Section A: Multiple-Choice Quiz \hfill \normalsize\normalfont(40 points --- 20 questions, 2 pts each)}
\textit{Select the single best answer. Questions marked [Basic] test core concepts; [Medium] test composite logic.}

\begin{enumerate}[label=\textbf{Q\arabic*.},leftmargin=1.8em,itemsep=0.6pt,topsep=0.2pt]

\item \textbf{[Basic]} At what precise moment does a standard Python \texttt{while condition:} evaluate its boolean test expression?
\begin{enumerate}[label=\Alph*.,leftmargin=1.8em]
  \item Only once when the program is initially compiled into bytecode before execution starts.
  \item Prior to every iteration cycle, including immediately before the very first iteration runs.
  \item Strictly after the entire loop body finishes executing, completely ignoring the first run.
  \item Only when an explicit \texttt{continue} statement forces an emergency re-evaluation.
\end{enumerate}

\item \textbf{[Basic]} What integer sequence is generated by evaluating \texttt{list(range(10, 2, -2))} in Python?
\begin{multicols}{4}
\begin{enumerate}[label=\Alph*.]
  \item \texttt{[10, 8, 6, 4, 2]}
  \item \texttt{[10, 8, 6, 4]}
  \item \texttt{[8, 6, 4, 2]}
  \item \texttt{[10, 7, 4, 2]}
\end{enumerate}
\end{multicols}

\item \textbf{[Basic]} Consider the following code snippet:\\
\texttt{total = 0}\\
\texttt{for x in range(3, 8): total += x}\\
What is the final value stored in variable \texttt{total} after loop termination?
\begin{multicols}{4}
\begin{enumerate}[label=\Alph*.]
  \item \texttt{22}
  \item \texttt{25}
  \item \texttt{30}
  \item \texttt{18}
\end{enumerate}
\end{multicols}

\item \textbf{[Basic]} In input-driven interactive loops, what is the defined role of a \textbf{sentinel value}?
\begin{enumerate}[label=\Alph*.,leftmargin=1.8em]
  \item A dedicated cryptographic token used to validate user network credentials.
  \item A predefined special value entered by the user to signal termination without being processed as valid data.
  \item A system counter that automatically tracks the number of CPU clock cycles elapsed.
  \item A required Python keyword that must precede every \texttt{while} loop statement.
\end{enumerate}

\item \textbf{[Basic]} What is the direct operational effect of executing a \texttt{break} statement within a loop?
\begin{enumerate}[label=\Alph*.,leftmargin=1.8em]
  \item It terminates the entire program immediately and shuts down the Python interpreter.
  \item It skips the remaining statements in the current iteration and jumps to the condition test for the next iteration.
  \item It immediately terminates execution of the innermost enclosing loop and jumps to the statement directly following it.
  \item It resets all loop variables to their initial values and restarts from the first iteration.
\end{enumerate}

\end{enumerate}

\newpage

% ==============================================================================
% PAGE 2: SECTION A (Q6 - Q13)
% ==============================================================================

\begin{enumerate}[label=\textbf{Q\arabic*.},leftmargin=1.8em,itemsep=0.8pt,topsep=0.3pt,resume]

\item \textbf{[Medium]} What is the final value of \texttt{res} evaluated by the following loop?\\
\texttt{res, m = 0, 1}\\
\texttt{while m <= 15:}\\
\texttt{\phantom{XX}if m \% 3 == 0: res += m}\\
\texttt{\phantom{XX}m += 2}
\begin{multicols}{4}
\begin{enumerate}[label=\Alph*.]
  \item \texttt{12}
  \item \texttt{24}
  \item \texttt{27}
  \item \texttt{45}
\end{enumerate}
\end{multicols}

\item \textbf{[Basic]} When searching for the minimum value across an arbitrary non-empty collection of positive floating-point numbers, how should the tracking variable \texttt{min\_val} be safely initialized?
\begin{enumerate}[label=\Alph*.,leftmargin=1.8em]
  \item Initialized to \texttt{0.0}, because 0 is the universal neutral number.
  \item Initialized to \texttt{-1.0}, to represent an unassigned state.
  \item Initialized to the first item of the sequence, or to positive infinity (\texttt{float('inf')}).
  \item Initialized to \texttt{None} and incremented by 1 inside each iteration.
\end{enumerate}

\item \textbf{[Basic]} What output is produced by the following character traversal loop?\\
\texttt{text, letters = "Py3.12-AI", 0}\\
\texttt{for ch in text:}\\
\texttt{\phantom{XX}if ch.isalpha(): letters += 1}\\
\texttt{print(letters)}
\begin{multicols}{4}
\begin{enumerate}[label=\Alph*.]
  \item \texttt{2}
  \item \texttt{4}
  \item \texttt{6}
  \item \texttt{9}
\end{enumerate}
\end{multicols}

\item \textbf{[Medium]} How many times is the statement \texttt{total += 1} executed in the following nested loops?\\
\texttt{total = 0}\\
\texttt{for a in range(1, 4):}\\
\texttt{\phantom{XX}for b in range(1, a + 1): total += 1}
\begin{multicols}{4}
\begin{enumerate}[label=\Alph*.]
  \item \texttt{6}
  \item \texttt{9}
  \item \texttt{12}
  \item \texttt{3}
\end{enumerate}
\end{multicols}

\item \textbf{[Basic]} In Python, under what condition does the \texttt{else} block attached to a \texttt{for} or \texttt{while} loop execute?
\begin{enumerate}[label=\Alph*.,leftmargin=1.8em]
  \item Only when the loop body encounters an explicit \texttt{break} statement.
  \item When the loop iterates to natural exhaustion or the condition evaluates to \texttt{False} without encountering \texttt{break}.
  \item At the start of every iteration before the loop condition is evaluated.
  \item Only when an unexpected runtime exception is raised inside the loop body.
\end{enumerate}

\item \textbf{[Basic]} What is the primary software engineering advantage of decomposing a program into modular functions?
\begin{enumerate}[label=\Alph*.,leftmargin=1.8em]
  \item It guarantees that the Python program will execute without allocating any memory on the heap.
  \item It promotes code reuse, simplifies debugging, reduces redundancy, and enhances system maintainability.
  \item It automatically converts dynamic Python scripts into compiled C++ machine binaries.
  \item It removes the need for conditional branches or loops throughout the entire codebase.
\end{enumerate}

\item \textbf{[Basic]} What is the fundamental distinction between \texttt{print()} and \texttt{return} inside a Python function?
\begin{enumerate}[label=\Alph*.,leftmargin=1.8em]
  \item \texttt{print()} transmits data back to the calling expression, while \texttt{return} outputs text to the screen console.
  \item \texttt{print()} displays text on the console without sending values back, while \texttt{return} exits the function and sends a computed value back to the caller.
  \item \texttt{return} can only be used with numbers, while \texttt{print()} can only be used with strings.
  \item There is no operational difference; both statements terminate the function immediately.
\end{enumerate}

\item \textbf{[Basic]} What happens to program control flow when a Python function invocation is encountered?
\begin{enumerate}[label=\Alph*.,leftmargin=1.8em]
  \item Control pauses at the call site, jumps to the called function's definition, executes its body, and returns to the caller.
  \item Both the caller and the called function continue executing simultaneously on separate CPU cores.
  \item The operating system terminates the calling process and launches an independent subprocess.
  \item Python queues the function to be executed only after the main script has completely ended.
\end{enumerate}

\end{enumerate}

\newpage

% ==============================================================================
% PAGE 3: SECTION A (Q14 - Q20) + SECTION B INTRO & PROBLEM B1
% ==============================================================================

\begin{enumerate}[label=\textbf{Q\arabic*.},leftmargin=1.8em,itemsep=0.7pt,topsep=0.3pt,resume]

\item \textbf{[Basic]} Which of the following function definition headers will cause Python to raise a \texttt{SyntaxError}?
\begin{enumerate}[label=\Alph*.,leftmargin=1.8em]
  \item \texttt{def format\_record(name, dept, role="Analyst"):}
  \item \texttt{def format\_record(name, dept="General", role="Analyst"):}
  \item \texttt{def format\_record(name="Guest", dept, role="Analyst"):}
  \item \texttt{def format\_record(name, dept, role):}
\end{enumerate}

\item \textbf{[Medium]} Given \texttt{def calc\_bill(u, base=12.0, rate=4.5): return u * rate + base}, what is evaluated by \texttt{calc\_bill(6, rate=5.0)}?
\begin{multicols}{4}
\begin{enumerate}[label=\Alph*.]
  \item \texttt{39.0}
  \item \texttt{42.0}
  \item \texttt{66.0}
  \item \texttt{57.0}
\end{enumerate}
\end{multicols}

\item \textbf{[Basic]} What is the \textbf{scope} of a local variable declared inside a Python function?
\begin{enumerate}[label=\Alph*.,leftmargin=1.8em]
  \item Accessible only within the function block where it was defined.
  \item Universally accessible across all modules and functions in the program.
  \item Retained in memory and accessible globally after the function finishes.
  \item Limited only to code executing inside \texttt{while} loops.
\end{enumerate}

\item \textbf{[Basic]} What occurs when a local variable inside a function shares the same name as a global variable?
\begin{enumerate}[label=\Alph*.,leftmargin=1.8em]
  \item Python immediately raises a compile-time \texttt{DuplicateNameError}.
  \item The local variable \textbf{shadows} the global variable within the local scope.
  \item The global variable permanently overwrites the local variable value.
  \item Python converts both variables into a combined two-element list.
\end{enumerate}

\item \textbf{[Medium]} What is printed by: \texttt{def double(n): n *= 2; return n} followed by \texttt{num = 24; double(num); print(num)}?
\begin{multicols}{4}
\begin{enumerate}[label=\Alph*.]
  \item \texttt{24}
  \item \texttt{48}
  \item \texttt{0}
  \item \texttt{None}
\end{enumerate}
\end{multicols}

\item \textbf{[Basic]} When a function returns two values using \texttt{return low, high}, what data structure is actually returned?
\begin{multicols}{4}
\begin{enumerate}[label=\Alph*.]
  \item \texttt{tuple}
  \item \texttt{list}
  \item \texttt{dict}
  \item \texttt{set}
\end{enumerate}
\end{multicols}

\item \textbf{[Medium]} In professional software engineering, what is the core rationale for designing functions with \textbf{high cohesion}?
\begin{enumerate}[label=\Alph*.,leftmargin=1.8em]
  \item Ensure each function performs a single, focused, and well-defined computational task.
  \item Permit functions to modify global variables directly without passing arguments.
  \item Force Python to compile dynamic scripts into optimized low-level machine instructions.
  \item Eliminate all loop structures in favor of recursive function invocations.
\end{enumerate}

\end{enumerate}

\vspace{-1.5mm}
\section*{Section B: Coding --- Fill-in Problems \hfill \normalsize\normalfont(30 points --- 2 problems, 10 blanks)}
\textit{Write missing Python expressions directly into designated answer slots. Difficulty: 24 pts Basic (80\%), 6 pts Medium (20\%).}

\subsection*{Problem B1: Daily Temperature Logging \& Extreme Detection (12 points --- 4 blanks, 3 pts each [All Basic])}
Complete the program below by providing expressions for blanks \textbf{[B1.1]} to \textbf{[B1.4]}.
The utility reads daily recorded temperatures (in Celsius, float) until the user enters the sentinel value \texttt{-999.0}, computes the cumulative sum, counts valid days, and prints a formatted summary.

\begin{lstlisting}
total_temp = 0.0
day_count = 0

# [B1.1] Priming read: prompt user for first temperature reading and cast to float:
temp = _________________________________________________________________  # Blank [B1.1] (3 pts - Basic)

# [B1.2] While temp is not equal to sentinel -999.0:
while __________________________________________________________________:  # Blank [B1.2] (3 pts - Basic)
    # [B1.3] Accumulate current temp into total_temp:
    total_temp = _______________________________________________________  # Blank [B1.3] (3 pts - Basic)
    day_count = day_count + 1
    temp = float(input("Enter next temperature (-999 to stop): "))

# [B1.4] Print final total formatted to 1 decimal place (e.g. 105.4):
print(f"Total: {______________________:.1f} C recorded across {day_count} days") # Blank [B1.4] (3 pts - Basic)
\end{lstlisting}

\begin{answerbox}[height=1.8cm]{Answer Slots for Problem B1}
\textbf{[B1.1]}: \dotfill \quad \textbf{[B1.2]}: \dotfill\\[1.5mm]
\textbf{[B1.3]}: \dotfill \quad \textbf{[B1.4]}: \dotfill
\end{answerbox}

\newpage

% ==============================================================================
% PAGE 4: SECTION B --- PROBLEM B2
% ==============================================================================

\subsection*{Problem B2: Modular Employee Sales \& Commission Processor (18 points --- 6 blanks, 3 pts each)}
\textit{Difficulty: Blanks [B2.1]--[B2.4] are Basic (12 pts); Blanks [B2.5]--[B2.6] are Medium (6 pts).}

Complete the modular employee commission processing utility below by providing expressions for blanks \textbf{[B2.1]} to \textbf{[B2.6]}.

\begin{lstlisting}
# Helper function to compute commission based on sales volume
# [B2.1] Basic: Define function header with default parameter rate=0.08:
def calc_commission(sales, _____________________________________________):  # Blank [B2.1] (3 pts - Basic)
    return sales * rate

sales_records = [1200.0, 4500.0, 8000.0, 2300.0, 6000.0]
total_commission = 0.0
max_sales = 0.0

# [B2.2] Basic: Traverse through each individual sales figure in sales_records:
for s in _______________________________________________________________:  # Blank [B2.2] (3 pts - Basic)
    # [B2.3] Basic: Call helper function calc_commission on s with default rate:
    comm = _____________________________________________________________  # Blank [B2.3] (3 pts - Basic)
    total_commission = total_commission + comm
    
    # [B2.4] Basic: Track the highest sales volume recorded so far:
    if s > max_sales:
        max_sales = ____________________________________________________  # Blank [B2.4] (3 pts - Basic)

# [B2.5] Medium: Call calc_commission on total_commission with executive bonus rate 0.15:
bonus_pool = ___________________________________________________________  # Blank [B2.5] (3 pts - Medium)

# [B2.6] Medium: Calculate average commission per rep safely avoiding ZeroDivisionError:
avg_commission = _______________________________________________________  # Blank [B2.6] (3 pts - Medium)

print(f"Bonus Pool: ${bonus_pool:.2f} | Max Sales: ${max_sales:.2f}")
\end{lstlisting}

\vspace{0.4mm}
\begin{answerbox}[height=1.8cm]{Answer Slots for Problem B2}
\textbf{[B2.1] (Basic):} \dotfill \quad \textbf{[B2.2] (Basic):} \dotfill\\[1.1mm]
\textbf{[B2.3] (Basic):} \dotfill \quad \textbf{[B2.4] (Basic):} \dotfill\\[1.1mm]
\textbf{[B2.5] (Medium):} \dotfill \quad \textbf{[B2.6] (Medium):} \dotfill
\end{answerbox}

\vspace{0.6mm}
% ==============================================================================
% SECTION C: IMMEDIATELY FOLLOWS SECTION B ON PAGE 4 (NO PAGEBREAK)
% ==============================================================================
\section*{Section C: Applied Full-Code Programming \& Design \hfill \normalsize\normalfont(30 points --- 3 problems)}
\textit{Difficulty: 24 pts Basic (80\%), 6 pts Medium (20\%). Problems C1 and C2 share the common logistics fleet scenario below with two distinct requirements. Problem C3 evaluates a sentinel loop and transaction auditing.}

\begin{tcolorbox}[colback=boxbg, colframe=secondaryteal, arc=1.2mm, boxrule=0.6pt, left=2.0mm, right=2.0mm, top=0.6mm, bottom=0.6mm]
\textbf{\small Problem Scenario for C1 \& C2: Logistics Delivery Fleet Fuel \& Trip Cost Analyzer}\\[0.2mm]
\footnotesize
A delivery logistics company tracks vehicle performance across $N$ delivery trips based on the following rules:
\begin{itemize}[itemsep=0.1pt,topsep=0.1pt]
  \item \textbf{Inputs per Trip:} Distance in km (\texttt{distance\_km}, float) and fuel rate in L/100 km (\texttt{liters\_per\_100km}, float).
  \item \textbf{Formulas per Trip:}
    $\text{Fuel Used (liters)} = \text{distance\_km} \times (\text{liters\_per\_100km} / 100.0)$; \quad
    $\text{Trip Cost (USD)} = \text{Fuel Used} \times 1.45$.
  \item Across all $N$ trips, the system tallies total distance (km), total fuel used (liters), total cost (USD), and identifies the maximum single-trip distance.
\end{itemize}
\end{tcolorbox}

\vspace{0.3mm}
\subsection*{Problem C1: Algorithmic Specification \& Pseudocode (8 points --- [All Basic])}
\textit{Requirement: Formulate structured \textbf{pseudocode} only (e.g., using \texttt{INPUT}, \texttt{COMPUTE} / \texttt{SET}, \texttt{FOR} / \texttt{WHILE}, \texttt{PRINT}). Do NOT write Python code.}

Write a complete algorithm in pseudocode that inputs the trip count $N$, iterates through all $N$ trips, computes each trip's fuel consumption and cost, and determines the \textbf{total distance}, \textbf{total fuel used}, \textbf{total cost}, and the \textbf{maximum single-trip distance}.

\begin{answerbox}[height=5.2cm]{Structured Pseudocode for Problem C1 (8 points --- Basic)}
\end{answerbox}

\newpage

% ==============================================================================
% PAGE 5: SECTION C --- PROBLEM C2 (MAXIMIZED CODING AREA)
% ==============================================================================

\subsection*{Problem C2: Modular Python Function Implementation (10 points --- [All Basic])}
\textit{Requirement: Write a clean, modular Python program implementing the delivery fleet engine. Define the calculation function and implement the main batch processing loop.}

Implement the program following these specifications:
\begin{enumerate}[label=\arabic*.,leftmargin=1.8em,itemsep=0.5pt]
  \item \textbf{Function \texttt{calculate\_trip\_stats(distance\_km, liters\_per\_100km, price\_per\_liter=1.45)}:}\\
    Compute and return both \texttt{fuel\_used} (float) and \texttt{cost} (float) as a tuple using the formulas from Problem Scenario C1 \& C2.
  \item \textbf{Batch Processing Loop:}
    Prompt the user for the number of trips $N$ (int). Use a loop to iterate $N$ times. In each iteration, prompt for \texttt{distance\_km} and \texttt{liters\_per\_100km}, call \texttt{calculate\_trip\_stats()}, and print each trip's itemized distance, fuel, and cost.
  \item \textbf{Accumulation \& Output Summary:}
    Tally total distance (km), total fuel used (L), and total cost (USD). After all trips are processed, print the summary:
\end{enumerate}

\begin{lstlisting}[basicstyle=\ttfamily\scriptsize]
Trip 1: Distance = 120.00 km | Fuel = 9.60 L | Cost = $13.92
Trip 2: Distance = 85.00 km | Fuel = 6.80 L | Cost = $9.86
--------------------------------------------------------
Total Trips: 2 | Total Distance: 205.00 km | Total Cost: $23.78
\end{lstlisting}

\vspace{0.8mm}
\begin{answerbox}[height=18.0cm]{Python Implementation for Problem C2 (10 points --- Basic Modular Functions \& Loops)}
\end{answerbox}

\newpage

% ==============================================================================
% PAGE 6: SECTION C --- PROBLEM C3 (SENTINEL ITERATION & TRANSACTION AUDITOR)
% ==============================================================================

\subsection*{Problem C3: Retail Store Daily Transaction Auditor \& High-Value Sale Analyzer (12 points)}
\textit{Difficulty Breakdown: 6 points Basic + 6 points Medium. Note: Requires a sentinel-controlled loop, modular high-value check helper function, and statistical performance tracking without complex branching.}

A store manager enters individual customer sales transaction amounts (in USD, positive float) until entering the sentinel value \texttt{-1.0}. Write a complete, concise Python program according to these specifications:

\textbf{Requirements:}
\begin{enumerate}[label=\arabic*.,leftmargin=1.8em,itemsep=0.5pt]
  \item \textbf{Modular Helper Function [3 pts --- Basic]:}\\
    Define a function \texttt{is\_high\_value(amount, threshold=100.0)} that returns \texttt{True} if $\text{amount} \ge \text{threshold}$, and \texttt{False} otherwise.
  \item \textbf{Sentinel Loop \& Data Aggregation [5 pts --- 3 Basic + 2 Medium]:}\\
    Use a \texttt{while} loop to prompt for transaction amounts until the sentinel \texttt{-1.0} is entered.\\
    Inside the loop, update:
    \begin{itemize}[itemsep=0.1pt]
      \item The total count of transactions.
      \item The cumulative sales revenue sum.
      \item The count of high-value transactions (by calling \texttt{is\_high\_value(amount)}).
      \item The highest single transaction amount recorded so far.
    \end{itemize}
  \item \textbf{Summary Report [4 pts --- Medium]:}\\
    If no transactions were entered (count is 0), print \texttt{"No transactions were recorded."} and terminate.\\
    Otherwise, compute:
    $\text{average} = \text{total revenue} / \text{count}$ \quad and \quad
    $\text{high-rate} = (\text{high count} / \text{count}) \times 100\%$.\\
    Display the formatted summary report (currency with 2 decimals, percentage with 1 decimal):
\end{enumerate}

\begin{lstlisting}[basicstyle=\ttfamily\scriptsize]
========================================================
DAILY TRANSACTION AUDIT SUMMARY
========================================================
Total Transactions: 5 | High-Value: 3 (60.0%)
Average Amount:     $142.50 | Highest Sale: $350.00
========================================================
\end{lstlisting}

\begin{answerbox}[height=8.8cm]{Python Implementation for Problem C3 (12 points --- 6 Basic + 6 Medium)}
\end{answerbox}

\vfill
\begin{center}
  \rule{0.5\textwidth}{0.4pt}\\[0.8mm]
  \textbf{\color{primaryblue}--- END OF EXAMINATION PAPER ---}
\end{center}

\end{document}
